
\documentclass[12pt]{article}

% Chinese support
\usepackage{xeCJK}
\setCJKmainfont{Microsoft YaHei}
\setCJKsansfont{Microsoft YaHei}
\setCJKmonofont{FangSong}

% Sans-serif font (fontspec for XeLaTeX)
\usepackage{fontspec}
\setmainfont{Arial}
\setsansfont{Arial}
\setmonofont{Courier New}

% Sans-serif math font
\usepackage[T1]{fontenc}
\usepackage{sfmath}
\usepackage{amssymb}

% Page margins
\usepackage[margin=0.1cm]{geometry}

% Section title sizes
\usepackage{titlesec}
\setcounter{secnumdepth}{4}
\titleformat{\section}{\fontsize{36}{43}\bfseries}{\thesection}{1em}{}
\titleformat{\subsection}{\fontsize{28}{34}\bfseries}{\thesubsection}{1em}{}
\titleformat{\subsubsection}{\fontsize{22}{27}\bfseries}{\thesubsubsection}{1em}{}
\titleformat{\paragraph}[runin]{\sffamily\fontsize{22}{27}\bfseries}{\theparagraph}{1em}{}

\begin{document}
\sffamily
\huge

我们的谱序列都要求
$$Z_\infty = \cap_n Z_n$$
$$B_\infty = \cup_n B_n$$


\section{filtered谱序列}
\subsection{filtered complex}
\paragraph{}filtered graded abelian group 的范畴：对象是decreasing filtration的具有n个指标的分次abelian group

\paragraph{}filtered complex是filtred abelian group范畴中的complex，带一个保filtration的微分d

\paragraph{} filtered complexd的同调上具有自然的filtratoin，给出证明

\paragraph{}associated graded 也构成一个complex

\paragraph{}我们假设完备性： $\cup F^nA =A$, $\cap F^nA=0$

\subsection{谱序列}
给定filtered complex A，
我们定义一个SSData的结构
\paragraph{}谱序列从$E_0$项开始，谱序列分次的个数为$n+1$

\paragraph{}定义$V$ 为A的associated graded

\paragraph{}定义$Z_n$为如下元素的像：
$x$的filtraton 为k，使得$dx$的filtration至少是n+k-1

\paragraph{}定义$B_n$为如下元素：
$x$的filtration为k，且形如$dy$，$y$的filtration至少k-n+1

\paragraph{}定义谱序列的微分为d自然诱导的映射，证明这个是良定义的

\paragraph{}证明如此的定义满足谱序列的公理

\paragraph{}证明按照我们的弱的收敛性的定义，这个谱序列收敛到A的同调

\end{document}
